\section{Động lực nghiên cứu}
\label{sec:dong_luc_nghien_cuu}

\noindent\textbf{Bối cảnh thực tiễn.} Té ngã là nguyên nhân hàng đầu gây tổn thương và tử vong ở người cao tuổi, với khoảng 37,3 triệu ca cần can thiệp y tế mỗi năm \cite{who2021falls}. Hậu quả không chỉ dừng ở chấn thương thể xác mà còn gây ra "cảm giác sợ ngã", tạo vòng xoắn bệnh lý khiến người bệnh e ngại vận động và tăng nguy cơ tái ngã \cite{dao2022thuongtich}. Đặc biệt, khi xảy ra khi người bệnh ở một mình mà không được phát hiện kịp thời, thời gian nằm trên sàn kéo dài có thể dẫn đến hạ thân nhiệt, hoại tử mô hoặc tử vong. Tại Việt Nam, quá trình già hóa dân số đang diễn ra nhanh chóng, đòi hỏi hệ thống giám sát tự động tại bệnh viện, viện dưỡng lão và hộ gia đình.

\noindent\textbf{Vấn đề chưa được giải quyết.} Các giải pháp hiện hành bộc lộ nhiều điểm yếu khi ứng dụng thực tế. Hệ thống dựa trên cảm biến đeo gây cảm giác vướng víu, bị giới hạn bởi dung lượng pin và tỷ lệ tuân thủ rất thấp ở người cao tuổi \cite{pham2025lightweight, khawam2025fall}. Trong khi đó, các hệ thống camera truyền thống phân tích trực tiếp hình ảnh RGB lại vi phạm nghiêm trọng quyền riêng tư \cite{ursul2024dynamic}. Thêm vào đó, các mô hình học sâu phân tích video hiện tại có cấu trúc cồng kềnh, không thể triển khai trên thiết bị biên Edge AI \cite{kurniadi2026human}.

\noindent\textbf{Hạn chế của các thuật toán hiện có.} Các phương pháp dựa trên ngưỡng cố định (\textit{threshold-based}) từ tín hiệu cảm biến có tỷ lệ báo động nhầm (\textit{false alarms}) cao, không phân biệt được cú ngã thực sự với các hoạt động sinh hoạt có gia tốc tương đồng \cite{benabdennour2026real}. Mạng LSTM gặp khó khăn trong việc nắm bắt phụ thuộc dài hạn do hiện tượng triệt tiêu đạo hàm, đồng thời thiếu cơ chế chú ý toàn cục để tập trung vào khoảnh khắc mất thăng bằng \cite{benabdennour2026real}. Việc xử lý trực tiếp trên luồng video thô cũng tạo độ trễ lớn, không đáp ứng yêu cầu của hệ thống cảnh báo thời gian thực.

\noindent\textbf{Đề xuất giải pháp.} Nhằm lấp đầy những khoảng trống trên, Khóa luận đề xuất hệ thống \textbf{HybridFallTransformer} -- giải quyết đồng thời bài toán về độ chính xác, tốc độ xử lý thời gian thực, bảo vệ quyền riêng tư và khả năng triển khai trên thiết bị biên. Hệ thống sử dụng YOLOv11n-Pose để trích xuất 17 điểm khớp COCO, chỉ xử lý vector tọa độ không gian mà không cần lưu trữ ảnh RGB thô \cite{khanam2024yolov11}. Dữ liệu khung xương được module PIFR biến đổi thành đặc trưng hình học -- động học \cite{kong2025pifr}, sau đó đưa vào bộ mã hóa Hybrid Transformer siêu nhẹ với cơ chế Temporal Self-Attention, khắc phục triệt để lỗi triệt tiêu đạo hàm của LSTM. Cơ chế cảnh báo qua Telegram được tích hợp để gửi thông báo tức thời đến người thân và nhân viên y tế, tối ưu hóa thời gian vàng trong cứu chữa.
